
% [Original opening paragraph retained as a comment — kept for potential future use]
% Sequential hypothesis testing requires two distinct decisions: \textit{when} to stop
% collecting data and \textit{what} to conclude once stopped. Conflating these steps,
% as in the HDI+ROPE algorithm, risks stopping on imprecise, unrepresentative posteriors
% and produces a systematic false-positive rate ($\sim\!6\%$ in our simulations, even
% when $\theta_{\rm true}=\theta_{\rm null}$).
% \citet{kruschke2015doing}'s Precision is the Goal (PitG) eliminates this systematic
% bias by decoupling the two steps: it stops only once the posterior HDI narrows below a
% pre-specified precision goal $\omega_{\rm goal}$, then applies the HDI+ROPE decision
% rule. In the fair-coin case, PitG achieves zero false positives, but at the cost of a
% high inconclusive rate of 63.3\%, because precision and decisiveness
% are not required to coincide.
\citet{kruschke2015doing} demonstrated that reliable sequential hypothesis testing
requires the stopping criterion to be decoupled from the decision criterion; PitG
achieves this, but at a high inconclusive rate --- up to $63.3\%$ in the fair-coin
case --- most critically when the null is true.

We have proposed Decisive Precision is the Goal (DPitG), a conservative variant of
PitG that joins the two criteria: data collection continues until the HDI is
simultaneously narrower than $\omega_{\rm goal}$ \textit{and} falls cleanly inside
or outside the ROPE\@.
This single change yields a $2.7\times$ gain in conclusiveness over PitG at a median
cost of only 4\% more samples ($\theta_{\rm true}=0.5$, $\omega_{\rm goal}=0.08$,
$\Delta_{\rm ROPE}=0.1$), while preserving a zero false-positive rate.
This advantage is robust: DPitG's conclusiveness varies by less than one percentage
point ($97.2\%$--$97.9\%$) across $\omega_{\rm goal}=0.06$--$0.10$, a range over
which PitG's conclusiveness collapses from $82.5\%$ to zero.
The gain is most pronounced near the ROPE boundaries, where correct decisions matter
most; once $\theta_{\rm true}$ moves well outside the ROPE, DPitG converges to PitG
and both stop at $N_{\rm goal}(\theta_{\rm true})$.

Unlike NHST (see Appendix~\ref{app:nhst}), which can only reject or fail to reject
the null, the HDI+ROPE decision rule enables positive \textit{acceptance} of the null
when the entire HDI falls within the ROPE --- a critical advantage for confirmatory
research areas such as bioequivalence, safety testing, and software regression testing;
DPitG ensures this verdict is reliably reached.

Prospective planning is straightforward: the closed-form
$N_{\rm goal}(\theta_{\rm null}, \omega_{\rm goal}) \propto
\theta_{\rm null}(1-\theta_{\rm null})/\omega_{\rm goal}^2$
provides the minimum budget required, with $N_{\rm max}$ set generously above this
baseline to ensure DPitG reaches its full conclusiveness potential.

The framework extends naturally to single-group continuous data
(Appendix~\ref{app:expected_stop}) and to two-group comparisons, both binomial
and continuous (Appendix~\ref{app:extensions}).\footnote{These extensions are
theoretical derivations; stopping-property validation beyond the single-group Bernoulli
case remains to be completed in a dedicated simulation study.}
All algorithms are implemented in openly available
code\footnote{\textcolor{red}{\textbf{[TODO: replace repo URL]}}
\href{https://github.com/elzurdo/precision-goal}{https://github.com/elzurdo/precision-goal}},
with an interactive online calculator for prospective sample-size planning available at
\href{https://r-w-t-y.streamlit.app}{https://r-w-t-y.streamlit.app}.

DPitG is the method of choice whenever a reliable verdict --- acceptance, rejection,
or inconclusive --- is required and the budget is sufficient to reach
$N_{\rm goal}(\theta_{\rm null})$; the closed-form $N_{\rm goal}$ makes this judgement
transparent before any data are collected.
